\documentclass[10pt]{article}
\usepackage[margin=20mm]{geometry}
\usepackage{amsmath,amssymb,mathtools,fontspec,unicode-math,xurl,hyperref}
\IfFontExistsTF{DejaVu Serif}{\setmainfont{DejaVu Serif}}{\setmainfont{Latin Modern Roman}}
\setmathfont{latinmodern-math.otf}
\setlength{\parindent}{0pt}\setlength{\parskip}{6pt}\sloppy
\title{Scientific initial review: M25-0001
Grid diagrams and Heegaard Floer invariants}
\author{}
\date{}
\begin{document}
\maketitle
Abstract: The reviewed object is the 55-page author manuscript arXiv:0910.0078v4, dated 1 July 2020, by Ciprian Manolescu, Peter S. Ozsváth and Dylan P. Thurston. Two defects affect steps used by the printed core construction or its finiteness argument. A third false general hyperbox statement has no demonstrated downstream core impact. Each finding has an explicit algebraic check or counterexample. These are findings about this manuscript version. The 2025 journal Version of Record was not acquired or compared; no claim about its current defects follows.\par
\section*{Object, scope, and version control}
The frozen journal identity is Annals of Mathematics 201 (2025), 1–78, DOI 10.4007/annals.2025.201.1.1 [S2]. The reviewed PDF and its bound author TeX are the earlier arXiv v4 [S1]. The publisher records a revision on 16 February 2021, after this v4. The arXiv history exposed in the version check lists v4 as its latest revision. A bounded search by the exact title and arXiv identifier with correction/erratum/error terms did not identify a separate correction. This does not establish the absence of corrections or intervening changes.\par
All 55 PDF pages, including references, were read as extracted text. The author TeX was inspected statically at the relevant definitions, formulas and proof passages. Sixteen source pages were also rendered and visually checked: 1, 3–8, 11, 23–25, 38, 42, and 44–46. Source-manifest bindings were checked by byte size and SHA-256. No author code was executed. No separately verified computational repository or supplement was available in the acquired input.\par
The claims are over \(\mathbb F_2\). Theorem 1.2 concerns completed minus and infinity groups and the plus group, with a sparse grid as input. Theorem 1.3 assumes a closed four-manifold with \(b_2^+>1\), a cut-link presentation, and a sparse grid. This review does not extend these claims to integer coefficients or arbitrary four-manifolds.\par
\section*{F01: the filtered projection is not defined on its printed domain}
\subsection*{Location and claim}
Equation reference: eq:printed. Section 2.2, pp. 7–8, the filtration defining \(\mathfrak A^-(G,s)\) and equation (4); bound file `hf+.tex`, lines 666, 669 and 674–679. The ordinary complex \(\mathbf A^-(G,s)\) is the part of \(C^-(G)\) with Alexander filtration \(A_i\leq s_i\). Its coefficient ring is \(\mathcal R=\mathbb F_2[[U_1,\ldots,U_n]]\), which has no negative \(U\) exponents. The intermediate complex has \[\mathfrak F_i(x)=-A_i(x),\quad
 \mathfrak F_i(U_{i,1})=-1,\quad \mathfrak F_i(V_{i,j})=1\ (j\geq2),\] and the other variables have filtration zero. The printed domain is \(\mathfrak F_i\leq s_i\). Put \(J=\{(i,j):1\leq i\leq\ell,\ 2\leq j\leq p_i\}\). On a monomial without \(Y\) factors, the asserted \(\mathcal R\)-linear projection is \[
 \widetilde P\left(\prod_{(i,j)\in J}V_{i,j}^{n_{i,j}}x\right)
 =
 \prod_{i=1}^{\ell} U_{i,1}^{A_i(x)-s_i-\sum_{j=2}^{p_i}n_{i,j}}
 \prod_{(i,j)\in J}U_{i,j}^{n_{i,j}}x .\]\par
\subsection*{An allowed sparse counterexample}
Use a \(4\) by \(4\) toroidal grid, with rows and columns \(0,1,2,3\). The following coordinates denote cells, with markings at their centres: \[\begin{array}{c|l}
\text{linked }O &(0,0),(2,2)\\
\text{free }O &(1,1),(3,3)\\
X &(0,2),(2,0).
\end{array}\] There is one \(O\) in every row and column, at most one \(X\) in each, and the rows and columns of the two free \(O\) markings contain no \(X\). The linked arcs form one rectangular component, an unknot. The \(X\) rows and columns are separated, including cyclically, by free rows and columns, so the grid is sparse.\par
Choose the state \(x\) formed by the four diagonal intersections \((0,0),(1,1),(2,2),(3,3)\) of the grid circles. These are vertices, distinct from marking centres. Put \(a=A_1(x)\in\mathbb Z\); the affine Alexander lattice for a knot is \(\mathbb Z\). Choose the permitted integer parameter \(s=|a|+1\). This argument does not need a particular absolute Alexander normalization. The single monomial \(x\) satisfies the boundedness condition (3), and \[\mathfrak F_1(x)=-a\leq s,\qquad
 \widetilde P(x)=U_{1,1}^{a-|a|-1}x,\qquad a-|a|-1\leq-1.\] It therefore lies in the printed domain but its purported image does not lie in \(C^-(G)\) over \(\mathcal R\), or in \(\mathbf A^-(G,s)\).\par
\subsection*{Counterevidence, local repair, and core impact}
Power-series completion does not admit negative exponents. The different Laurent coefficient ring used later for the infinity theory does not rescue this minus map. Sparsity does not exclude the example. No stated reindexing of \(x\) or \(s\) changes these explicit definitions.\par
Changing the cutoff to \(\mathfrak F_i\leq -s_i\) is a plausible local repair. For \(U_{1,1}^r\prod_{j\geq2}V_{1,j}^{n_j}x\), set \(N=\sum_jn_j\). The changed inequality would give \[-a-r+N\leq-s\ \Longrightarrow\ r+a-s-N\geq0.\] The image then has Alexander filtration \(a-(r+a-s-N)-N=s-r\leq s\). Additional nonnegative \(U\) powers only decrease that target filtration. This proves this local well-definedness calculation; it does not establish the quasi-isomorphism or globally repair the subsequent conventions.\par
Equation (9), p. 11, repeats the exponent in the projection-inclusion map \(I_s^{\vec M}\). Taking the whole knot as sublink with its given orientation reproduces the obstruction. Equation (29), p. 23, composes this map into \(\Phi_s^{\vec M}\). These maps enter the surgery differential on p. 24, Theorem 2.10 on p. 25, and the construction used for Theorem 1.2 on pp. 45–46. Thus a printed core construction is undefined on an allowed input. Classification: Major in v4. The computability theorem itself is not disproved; a consistently propagated sign correction may repair the construction.\par
\section*{F02: local hyperbox relations do not imply the asserted total differential}
\subsection*{Location and claim}
Section 2.1, p. 5, equation (1) and the subsequent total-complex paragraph; p. 6, the induced map of those total complexes. Bound TeX: lines 514–528 and 553. Relations (1) are imposed for displacements \(\varepsilon\in\{0,1\}^n\). The text then claims that the sum of local maps is a differential on \[C_*=\bigoplus_{\varepsilon\in E(\mathbf d)}
 C^\varepsilon_{*+\|\varepsilon\|}\] for a general hyperbox, before compression.\par
\subsection*{Complete finite counterexample}
Let \(n=1\), \(d_1=2\), and \(C^0=C^1=C^2=\mathbb F_2\), each concentrated in internal degree zero. All vertex differentials vanish, and both consecutive edge maps are identities. Their internal degree is zero, as required. Relations (1) for displacement zero are \(d^2=0\) and those for displacement one are \(df+fd=0\); all hold. No condition is imposed for displacement two.\par
On the total space let \(e_j\) be the vector at vertex \(j\), of total degree \(-j\). Then \[De_0=e_1,\quad De_1=e_2,\quad De_2=0,\quad D^2e_0=e_2\ne0.\] With columns representing inputs, the independent matrix calculation over \(\mathbb F_2\) is \[D=\begin{pmatrix}0&0&0\\1&0&0\\0&1&0\end{pmatrix},
 \qquad
 D^2=\begin{pmatrix}0&0&0\\0&0&0\\1&0&0\end{pmatrix}.\] There is no convergence or sign ambiguity.\par
\subsection*{Counterevidence and limited impact}
In a unit hypercube all endpoint displacements belong to \(\{0,1\}^n\), so (1) supplies every component of \(D^2=0\). Compression of this example gives the two-vertex cube \(C^0\xrightarrow{\operatorname{id}}C^2\), whose total differential squares to zero. Restricting the claim to hypercubes, or compressing first, removes the obstruction.\par
The descent construction and pp. 42, 45 explicitly use compression. The surgery complex has the unit-cube indexing by sublinks. No core dependence on the false totalization of an uncompressed general hyperbox was established. Classification: Minor. The p. 5 statement and related p. 6 assertion need qualification; this example does not invalidate the later surgery formula.\par
\section*{F03: a global finite-dimensional truncation cannot preserve the minus tower}
\subsection*{Location and claim}
Page 46, first paragraph in the proof of Theorem 1.2; bound TeX lines 2476–2477. After acknowledging the completed minus complex is infinite dimensional, the proof claims it can be replaced by a quasi-isomorphic complex finite dimensional over \(\mathbb F_2\), using vertical and horizontal truncation in manuscript reference [14, Section 10]. This supplies the final finiteness step.\par
\subsection*{Counterexample and complete dimension obstruction}
Take the allowed sparse unknot grid in F01 and framing \(+1\). Surgery gives \(S^3\), with its unique \(\operatorname{Spin}^c\) structure. The genus-zero, one-generator minus Floer complex of \(S^3\) has zero differential: the symmetric product of degree zero is a point and there is no index-one disk contributing a differential. With the completion explicitly used by Theorem 1.2, \[\mathbf{HF}^{-}(S^3)\cong\mathbb F_2[[U]].\] For every \(N\geq0\), the elements \(1,U,\ldots,U^N\) are linearly independent by uniqueness of formal power-series coefficients. The homology thus has infinite \(\mathbb F_2\)-dimension.\par
A complex of finite total dimension \(m\) has homology \(\ker d/\operatorname{im}d\) of dimension at most \(m\). A quasi-isomorphism preserves homology. Consequently no such finite complex can be quasi-isomorphic to this completed minus complex. This analytic argument, not a finite experiment, establishes the obstruction.\par
\subsection*{Alternative interpretations and consequence}
A finite presentation or finite rank over \(\mathbb F_2[[U]]\) is compatible with the example. A finite grading window or a specified-degree calculation can also be meaningful. These are different from the unqualified global assertion of finite dimension over \(\mathbb F_2\) on p. 46. Surrounding discussion of ranks motivates checking a narrower intended reading, but that restriction is not stated in this step.\par
Infinite groups can have finite algorithmic descriptions, so computability itself is not refuted. A repair must specify the output representation, retained tower information, and any degree restriction, and show that the chosen truncation recovers that object. This review does not certify that replacement argument or re-prove the external truncation theorem. Classification: Major in v4: the printed core finiteness justification fails.\par
\section*{Positive checks and coverage limitations}
For the last numerical inference in Proposition 3.10, the inequality \(I\geq(\epsilon_1-1)/2\), with \(\epsilon_1\in\{0,1\}\) and \(I\in\mathbb Z\), does imply \(I\geq0\). This checks the final inference, not every preceding geometric step.\par
The use of Theorem 4.4 in Lemma 4.8 has the required degree range: for every \(k\geq2\), \(1-k<0\) and \(1-k<2-k\), hence \(1-k<\min(0,2-k)\). In the cocycle calculation on pp. 44–45, substituting the induction relations produces associative triple products occurring in pairs, which cancel in characteristic two.\par
The text was read for definitions, quantifiers, coefficients, hypotheses, boundary values and dependency directions. This did not amount to independent reconstruction of all holomorphic-disk counts, transversality arguments, sparse-domain case analyses, or external surgery results. No further supported issue is asserted. This is not a certificate of correctness for every proof.\par
\section*{Conclusion}
ID | Assessment | Established consequence\par
F01 | Major, v4 | A printed core projection leaves its coefficient ring on an allowed sparse input.\par
F02 | Minor, v4 | Local hyperbox relations do not ensure the uncompressed total differential; later compression limits impact.\par
F03 | Major, v4 | A finite-total-dimensional \(\mathbb F_2\) complex cannot retain the completed minus homology of \(S^3\).\par
The mathematical clarifications needed are a consistent filtration sign convention; an appropriately restricted totalization statement; and a precise finite representation or degreewise algorithm in the last step of Theorem 1.2. Compare the exact journal text and later corrections before assigning any finding to the 2025 publication. No editorial decision or author contact is part of this initial review.\par
\section*{Supporting checks}
The reviewer-authored `independent\_checks.py` uses Python standard-library arithmetic and JSON. Its actual output records the matrix square, grid incidence/sparsity checks, and 25 F01 instances with \(a=-12,\ldots,12\). It also checks 9,465 admitted monomials for the proposed local cutoff repair, with \(a,s=-5,\ldots,5\), \(N=0,\ldots,6\), and \(r=0,\ldots,14\). No negative output exponent or target-filtration violation occurs. The universal proof is in F01; the scan is corroboration only. The degree inequality was also evaluated for \(k=2,\ldots,14\).\par
The evidence directory includes the extraction, rendered source pages and logs, independent script and actual output, and coverage, findings and provenance records.\par
4ccd1086d715e547c243ca446da942e472a77cede7830b44322397f34a2cd503\par
Exact report-source, compiler, dependency and output-PDF hashes are in the external build and provenance records; a self-referential PDF hash is deliberately not embedded here.\par
\section*{Execution record and remaining limits}
Scientific analysis, independent calculations, static source checks, report creation and visual checks were performed by the assistant. No backend equivalence was inferred, no independent routing evidence established a mismatch, and no model switch was requested or performed.\par
The existing native `/usr/bin/pdflatex` was independently verified. The user-authorized portable Tectonic 0.17.0 binary was hash verified. Author TeX was inspected, not executed. No software installation, author-code execution, shell escape, or author communication was used.\par
A direct PDF extraction tool failed; installed Poppler `pdftotext` succeeded. A console read encountered a Windows character-encoding error and was rerun read-only with UTF-8 output. An arXiv HTML fetch for the external surgery paper timed out, so that attempt did not provide its complete proof. These failures are not evidence of a scientific defect. Human adjudication and full-text comparison with the journal publication remain outside this completed, version-limited initial review.\par
S3 C. Manolescu, P. S. Ozsváth and D. P. Thurston, Grid diagrams and Heegaard Floer invariants, arXiv:0910.0078v4, 1 July 2020, 55 pp. Reviewed PDF and bound author source. https://arxiv.org/abs/0910.0078. Annals of Mathematics, publisher record, 201 (2025), 1–78. DOI: 10.4007/annals.2025.201.1.1. Metadata checked; journal full text not reviewed. https://annals.math.princeton.edu/2025/201-1/p01. C. Manolescu and P. S. Ozsváth, Heegaard Floer homology and integer surgeries on links, arXiv:1011.1317; v6 metadata dated 16 July 2024. Manuscript reference [14]. Metadata and dependency use checked; complete external proof not independently verified. https://arxiv.org/abs/1011.1317.\par
\end{document}
